\section{Complete recognition from a checked two-meridian derivation}
\label{sec:two-meridian}

\subsection{The delivered idea and the maintained topology adapter}
Commit \code{ea9cf443c} supplied a particularly useful consequence of the
traceless-representation theorem: recognition is polynomial when every
Wirtinger generator can be derived from two original meridians by crossing
conjugations. The maintained implementation now checks this condition directly
on a validated planar diagram. This produces actual two-sided knot verdicts
on a restricted class. Exhaustion of the search says nothing about knot type.

The new module is \code{fastunknot.two\_meridian}.

Enable the stage with \path{--two-meridian} or Python's
\path{use_two_meridian=True}. It runs after the existing invariant filters and
before optional group simplification or the remaining exact fallback. The
ordinary defaults are unchanged. A successful negative test of representation
existence returns \code{UNKNOT}; a successful existence test returns
\code{KNOTTED}. Failure to obtain the required derivation, a local timeout or
local work exhaustion returns \code{INCONCLUSIVE} and resumes the pipeline.

The adapter revalidates the spherical, one-component PD instead of trusting a
presentation-level declaration of meridian provenance. At each crossing,
union the two overpassing edge ports. The resulting classes are the actual
Wirtinger arcs. The oriented underpassing ports and crossing sign give
\[
 x_v=x_o^{\epsilon}x_u x_o^{-\epsilon},\qquad\epsilon\in\{1,-1\}.
\]
Both directions are retained: knowing $x_o,x_u$ derives $x_v$, and knowing
$x_o,x_v$ derives $x_u$. The adapter hashes the canonical validated PD and
binds every certificate to that input. A connected-sum factor uses its own
PD and its own meridians; the original source braid is not substituted for a
factor presentation.

\subsection{A seed derivation is an elimination certificate}
A seed set consists of one or two \emph{original} Wirtinger arcs. A queue
propagates the crossing rules whose dependencies have become known. A rule
with identical over and source arcs has one distinct dependency, not two.
Every accepted step records a previously unknown target, the original
crossing index and the direction. Its two dependencies have already been
derived, so the trace is acyclic and has one step per nonseed arc.

For a fixed seed set this computes the least closed set of known arcs.
Induction shows that every queued step belongs to every rule-closed superset
of the seeds. Conversely, when the queue is empty, every rule with known
dependencies has already fired. Therefore queue failure cannot be repaired
by choosing a different order of rules with those same seeds. The search
tries all single seeds and then all unordered pairs, subject to its caps.
This recognizes the diagram-specific two-seed class exactly when uncapped.
It does not search all diagrams of the same knot.

Replay uses the recorded rules to define a word $w_a$ in the seed meridians
for every arc $a$. The chosen crossing relation is a valid Tietze elimination
of its new target. Substituting the definitions into \emph{every original
crossing relation}, including unused ones, preserves the original group
presentation. Equivalently, assigning values to the seeds extends uniquely
along the trace, and is a representation of the knot group exactly when all
substituted original relations hold. No omitted residual relation can create
a spurious representation. An arbitrary two-generator presentation obtained
by a separate group search does not suffice: its generators need not be
meridians, even if their abelianized exponents are $\pm1$.

The delivered prototype also minimizes formal degrees for fixed seeds.
The maintained two-seed path needs only a valid derivation and uses ordinary
queue closure. It does not implement that degree optimization or the report's
general multivariate real-feasibility compiler.

\subsection{Why two traceless meridians reduce to integer arithmetic}
Kronheimer--Mrowka's Corollary 7.17 states that every nontrivial knot admits an
irreducible $\mathrm{SU}(2)$ representation sending a chosen meridian to $i$
\cite{km-traceless-primary}. All positively oriented Wirtinger meridians are
conjugate; their images therefore have trace zero. Conversely, the infinite
cyclic group of the unknot has only abelian images. For a checked two-seed
presentation, nontriviality is thus equivalent to existence of noncommuting
unit pure imaginary quaternion values $A,B$ satisfying the substituted
relations. If the seeds commute, every derived value commutes, so such a
representation is not irreducible.

Set $T=AB$. Since $A^2=B^2=-1$ and $ATA^{-1}=T^{-1}$, all word values have
normal form
\[
 (-1)^eT^kA^s,\qquad e,s\in\{0,1\},\quad k\in\mathbb Z.
\]
The seed triples are $A=(0,0,1)$ and $B=(1,-1,1)$. Direct multiplication gives
\begin{align*}
 (e,k,s)(f,\ell,t)&=(e+f+st\bmod2,\ k+(-1)^s\ell,\ s+t\bmod2),\\
 (e,k,s)^{-1}&=(e+s\bmod2,\ (-1)^{s+1}k,\ s).
\end{align*}
These operations evaluate all derivation words without expanding them.
The checker forms each original relation as the predicted outgoing meridian
times the inverse of the actual outgoing value.

For noncommuting $A,B$, write $T=\cos\phi+J\sin\phi$ with
$0<\phi<\pi$. Every such phase occurs, for example with
$A=i$ and $B=-\cos\phi\,i+\sin\phi\,j$. Terms with $s=1$ lie in the
orthogonal quaternion plane and cannot equal the identity. For a relation
with $s=0$, the remaining condition is
\[
 k_j\phi\equiv e_j\pi\pmod{2\pi}.
\]
A relation with $k_j=0,e_j=1$ is immediately inconsistent. Otherwise put
$d=\gcd_j|k_j|$. If $d=0$, all relations are identities and any interior
phase, such as $\pi/2$, supplies a noncommuting representation.

If $d>0$, at least one $k_j/d$ is odd. Its central bit fixes a parity $e$.
The equations are compatible exactly when
\[
 e_j\equiv(k_j/d)e\pmod2\quad\hbox{for every }j.
\]
To see necessity, a B\'ezout combination of the integers $k_j/d$ shows that
$d\phi/\pi$ is an integer $q$, and each equation fixes
$(k_j/d)q\equiv e_j\pmod2$. An odd quotient determines $q\bmod2$.
Conversely, the displayed compatibility makes every integer $q\equiv e\pmod2$
a simultaneous solution with $\phi=q\pi/d$.
The strict interior condition requires $1\le q\le d-1$, so there are
\[
 \left\lfloor\frac{d-1+e}{2}\right\rfloor
\]
allowed phases. A positive count supplies the compact witness $q=1$ for odd
parity or $q=2$ for even parity. A zero count leaves only commuting endpoint
phases and certifies the unknot. Floating-point trigonometry and a general
real-algebraic solver are unnecessary.

For one seed, every crossing derivation is formally the same meridian: a
conjugate of $A$ by $A^{\pm1}$ is $A$. All original relations become identities,
the group is cyclic, and the diagram is an unknot. The zero-crossing diagram
is handled directly.

\subsection{Completeness, complexity and certificate scope}
The uncapped standalone algorithm is complete on diagrams admitting a derivation from
at most two seeds. There are $O(n^2)$ seed sets; each closure uses $O(n)$
incidence and queue operations. Indices have $O(\log(n+2))$ bits. Hence seed
search costs $O(n^3\log(n+2))$ bit operations with an elementary model.

For a successful derivation, the maximum absolute exponent grows by at most
a fixed factor plus a constant at each step, so it is at most $3^{O(n)}$.
Every stored integer consequently has $O(n)$ bits. There are $O(n)$ word
operations and relations. Even allowing $O(n)$ long-division steps at
$O(n^2)$ cost for each of the $O(n)$ gcd computations gives the conservative
complete bound
\[
 O(n^4\log(n+2))
\]
on this class, including PD preparation, substitution, arithmetic and replay.
This deliberately loose bound does not depend on expanding a word of length
$3^{O(n)}$. The trace and relation triples occupy $O(n^2)$ bits. More efficient
integer arithmetic can improve this accounting but is not needed for
polynomiality. This bound is for the standalone stage. The optional pipeline
placement retains earlier invariant filters, whose costs must be added;
an expensive Jones computation before this stage is not covered by this
restricted-class polynomial bound.

This theorem does not cover arbitrary input knots, nor prove a cheap
conversion of every diagram into the two-seed class. In particular, seed
failure is not evidence of nontriviality. The general recognition fallback
still has an exponential worst-case bound. The maintained faithful Jones
query has a general $2^{O(\sqrt n)}$ bound, with fixed polynomial factors
absorbed, but equality of Jones polynomials does not finish recognition.
The new two-seed theorem supplies another complete restricted class, not the
missing general quasi-polynomial theorem.

A certificate records the version, input hash, arc count, seed indices,
derivation trace, every relation triple, the gcd/parity calculation and the
verdict. The public verifier reconstructs the signed presentation from the
PD and replays the trace without rerunning seed search. It recomputes all
relation triples and the final answer and compares the complete strict
schema, including integer types. Producer and verifier share the PD parser
and integer arithmetic; this is search-independent replay, not an independent
implementation of topology or quaternion multiplication. Independent exact
quaternion and whole-cube comparisons test that shared arithmetic separately.

The local defaults are 50 milliseconds, two million weighted work units and
10,000 attempted seed sets. Preparation, failed attempts, successful
arithmetic and the mandatory replay share the same allowance; replay does
not receive a fresh budget. The counter charges integer bit lengths but is a
cooperative work metric, not an exact bit-operation count or a hard memory
bound. Deadline and caller cancellation checks occur throughout the search
and replay. Exhaustion before verified publication emits no certificate.
A global pipeline limit retains the existing \code{UNKNOWN} semantics.
Python callers can set all three local caps to \code{None} for exhaustive
restricted-class recognition; the default finite caps are a practical probe.

\subsection{Validation and full-pipeline measurements}
Eight focused tests cover 300 rational-quaternion product/inverse comparisons,
120 small phase systems compared with explicit rational phases, signed PD
relations compared with the existing full Wirtinger builder, and 80 seeded
small classical diagrams compared with an independent expanded Khovanov cube.
For every seed set of those small inputs, a separate fixed-point closure
agrees with the queue algorithm. Mirrors, rotated underpassing ports and
permuted crossing rows exercise orientation handling. Replay is tested with
seed search disabled, forged schema fields and derivations, a mismatched PD,
large integer JSON encoding, exact work boundaries, caller cancellation,
local fallback and CLI dispatch. The connected-sum test checks the actual
interlacement factor and rejects replay against the whole input. All 756
integrated tests pass with Regina installed in 174.738 seconds; the eight
focused methods pass in 1.131 seconds.

The benchmark fixes the previous recognizer at \code{7c45c41bc}, loads it
under the unchanged package dependencies, and runs it twice as an A/A control.
The other arms use the maintained recognizer with the new stage disabled or
enabled. Source hashes are frozen before timing and verified afterward.
All four arms revalidate a fresh PD and run the whole pipeline, including
filters, certificate verification and fallback. Import time and construction
of the previously reduced hard-input corpus are excluded. The corpus consists
of the 15 existing compressed-group hard-unknot cases and five named controls:
trefoil, figure-eight, Hard 8, Conway and Kinoshita--Terasaka.

Each input is tested under the ordinary pipeline, with a three-second global
allowance, and the stronger compressed-group configuration, with 15 seconds
for group search, 20 million group work units and an 18-second global
allowance. Both configurations retain all existing filters and allow 50,000
Khovanov objects. There is one excluded warmup per input, configuration and
arm, then five measured rounds with shuffled arm order. The optional stage
uses its ordinary 50-millisecond local allowance. This gives 800 measured
whole-recognition calls and 160 excluded warmup attempts. A separate one-shot
standalone probe on each input omits the local time cap while retaining the
work and attempt caps; those are discovery observations, not repeated timing
comparisons. All known labels predate this recognizer.

Raw samples, completion status, local work counts and deduplicated certificate
contents are retained in:
\begin{center}\small
\path{fast/results/two_meridian_20261008.json}
\end{center}
Completed-time medians exclude censored runs. A paired ratio is formed only
when both queries complete; it is the median of per-round ratios, not a ratio
of medians. The following table displays timings only for arms completing
all five measured queries and ratios only for five complete pairs.

\begin{center}\small
\begin{tabular}{@{}lrrrrrrr@{}}
\toprule
Input & $n$ & Ord. off & Ord. on & Ratio & Group off & Group on & Ratio \\
\midrule
Survivor 00 & 16 & 12.45 & 3.14 & 3.84 & 18.05 & 3.09 & 5.85 \\
Survivor 01 & 15 & 13.90 & 15.40 & 0.90 & 16.49 & 17.71 & 0.91 \\
Survivor 02 & 15 & 5.05 & 2.19 & 2.26 & 16.46 & 2.02 & 8.18 \\
Survivor 03 & 13 & 3.97 & 1.79 & 2.22 & 11.83 & 1.84 & 6.45 \\
Mirror 03 & 13 & 4.05 & 1.83 & 2.33 & 25.49 & 1.80 & 14.15 \\
Survivor 04 & 13 & 3.49 & 1.76 & 1.97 & 10.37 & 1.84 & 5.55 \\
Survivor 05 & 13 & 4.38 & 1.84 & 2.40 & 19.91 & 1.86 & 10.79 \\
Survivor 06 & 12 & 3.98 & 1.84 & 2.16 & 7.89 & 1.85 & 4.27 \\
Survivor 07 & 12 & 3.48 & 1.88 & 1.85 & 22.84 & 1.94 & 11.55 \\
Survivor 08 & 11 & 3.13 & 1.51 & 2.06 & 7.79 & 1.52 & 5.14 \\
Mirror 08 & 11 & 3.16 & 1.53 & 2.04 & 20.88 & 1.61 & 13.08 \\
Survivor 09 & 11 & 3.08 & 1.52 & 2.00 & 7.80 & 1.53 & 5.17 \\
Survivor 10 & 11 & 2.94 & 1.53 & 1.94 & 7.10 & 1.60 & 4.45 \\
Survivor 11 & 11 & 3.37 & 1.63 & 2.07 & 7.68 & 1.66 & 4.61 \\
Gordian & 141 & limit & limit & -- & 2091.80 & 2113.05 & 0.99 \\
Trefoil & 3 & 0.06 & 0.06 & 1.00 & 0.06 & 0.06 & 1.00 \\
Figure eight & 4 & 0.06 & 0.06 & 1.00 & 0.06 & 0.06 & 1.01 \\
Hard 8 & 8 & 0.68 & 1.02 & 0.67 & 2.56 & 1.05 & 2.41 \\
Conway & 11 & 1.04 & 1.01 & 1.03 & 1.02 & 1.01 & 1.00 \\
KT & 11 & 1.05 & 1.06 & 1.00 & 1.05 & 1.05 & 1.00 \\
\bottomrule
\end{tabular}
\end{center}
Median milliseconds for whole recognition, including PD validation and replay. Ratios are medians of paired off/on times; limits are censored.

The standalone stage resolves 13 of the 15 established hard-unknot inputs.
It exhausts all 120 seed sets on survivor 01, reaching at most seven of its
15 arcs. Gordian reaches the two-million-unit work cap after 4,609 attempts;
that capped failure does not determine whether its diagram belongs to the
class. Trefoil, figure-eight and Hard 8 are also conclusive; Conway and
Kinoshita--Terasaka exhaust their 66 seed sets and remain inconclusive.
Their ordinary pipeline verdicts come from earlier Jones obstructions.

For the 13 successful hard-corpus cases, ordinary full-recognition paired
speedups range from $1.85$ to $3.84$, with a median of the 13 per-case ratios
of $2.07$. Against the compressed-group configuration, the corresponding
range is $4.27$--$14.15$ and median $5.85$. For example, mirror 03 changes
from 25.49 to 1.80 milliseconds in the group configuration; the ordinary
configuration changes from 4.05 to 1.83 milliseconds. These are measured
whole-query improvements with replay included, on an existing finite corpus,
not an asymptotic claim or an independent random-knot sample.

There are material counterexamples to unconditional dispatch. Hard 8 changes
from 0.68 to 1.02 milliseconds in the ordinary pipeline: the original late
path is cheaper than constructing this certificate. Its paired off/on ratio
is $0.666$, about a 50\% slowdown, although the same probe improves the
optional group configuration by a factor of $2.41$. The unsuccessful search
on survivor 01 gives ordinary/group ratios $0.903$ and $0.914$. Gordian's
50-millisecond probe times out in both configurations. All 20 measured
ordinary Gordian queries exhaust the global allowance and remain censored;
all group-configured Gordian queries complete using the original group
certificate, with disabled/enabled medians 2.092/2.113 seconds and paired
ratio $0.987$. No timeout is reported as a successful three-second decision.

Overall, 780 of the 800 measured queries complete, every completed verdict
agrees with the prior label, and 20 are censored. Ordinary A/A ratios range
from $0.965$ to $1.027$ on completed pairs; group A/A ratios range from
$0.946$ to $1.052$. Comparing the old recognizer to the current disabled arm
also retains an ordinary survivor-00 outlier of $0.882$, so small timing
changes are not interpreted as improvements from argument validation or
other incidental work. The large successful-stage improvements coexist
with the retained regressions and limited corpus size. The implementation
therefore remains opt-in. Choosing when to attempt it earlier, or when to
skip it in favor of an already cheap path, requires a separately measured
adaptive policy.
